\appendixsection{Kafka-топики и поток сообщений}
\label{app:kafka}

\subsection{Основные топики платформы}

Kafka используется как транспорт асинхронного обмена между фасадом, планировщиком,
sender-сервисами и сервисом истории. Актуальная карта топиков приведена в
таблице~\ref{tab:appendix-kafka-topics}.

\begin{table}[h]
    \centering
    \footnotesize
    \renewcommand{\arraystretch}{1.15}
    \caption{Назначение Kafka-топиков}
    \label{tab:appendix-kafka-topics}
    \begin{tabular}{|L{0.24\textwidth}|L{0.20\textwidth}|L{0.20\textwidth}|L{0.24\textwidth}|}
        \hline
        \textbf{Топик} & \textbf{Продюсер} & \textbf{Консьюмер} & \textbf{Основные типы событий} \\
        \hline
        notification.facade.events & Notification Facade & Mail Sender, SMS Sender & EventCreated \\
        \hline
        notification.facade.dispatches & Notification Facade & Нет выделенного бизнес-потребителя в текущем коде & События по dispatch \\
        \hline
        notification.mail.dispatches & Notification Facade, Scheduler Delivery & Mail Sender & MailDispatchRequested \\
        \hline
        notification.mail.dispatches.scheduled & Notification Facade & Scheduler Delivery & MailDispatchRequested с будущим planned\_send\_at \\
        \hline
        notification.sms.dispatches & Notification Facade & SMS Sender & SmsDispatchRequested \\
        \hline
        notification.mail.delivery-statuses & Mail Sender & History Writer & MailDeliveryStatusChanged \\
        \hline
        notification.sms.delivery-statuses & SMS Sender & History Writer & SmsDeliveryStatusChanged \\
        \hline
    \end{tabular}
\end{table}

\subsection{Формат интеграционного сообщения}

Сообщения, публикуемые из outbox, имеют единый логический конверт:
\begin{enumerate}
    \item outboxId --- идентификатор записи выходящего журнала;
    \item aggregateType --- тип агрегата, например notification\_event, mail\_dispatch, sms\_delivery;
    \item aggregateId --- идентификатор агрегата;
    \item eventType --- прикладной тип события;
    \item payload --- полезная нагрузка бизнес-события;
    \item headers --- служебные заголовки, включая message\_id и дополнительные идентификаторы;
    \item createdAt --- время создания outbox-записи.
\end{enumerate}

Унификация формата позволяет sender-сервисам и history-writer использовать
одинаковую схему десериализации и входящего хранения, а также делает трассировку
сообщений устойчивой к расширению числа событий.

\subsection{Поток EventCreated}

Поток событий создания уведомления организован следующим образом:
\begin{enumerate}
    \item фасад создаёт запись notification\_event;
    \item в outbox добавляется сообщение EventCreated типа notification\_event;
    \item relay фасада публикует сообщение в notification.facade.events;
    \item Mail Sender и SMS Sender читают сообщение, сохраняют его во входящий журнал и далее строят собственные записи доставки только для своего канала.
\end{enumerate}

Такой подход позволяет sender-сервисам реагировать на факт создания события, не
завися от немедленного вызова TriggerDispatch.

\subsection{Поток TriggerDispatch}

После вызова TriggerDispatch фасад создаёт запись dispatch, snapshot аудитории
и канальную outbox-команду:
\begin{enumerate}
    \item для email --- MailDispatchRequested;
    \item для SMS --- SmsDispatchRequested.
\end{enumerate}

Если email-запуск отложен, OutboxRelayService направляет сообщение в
notification.mail.dispatches.scheduled; если момент отправки уже наступил, оно
сразу публикуется в notification.mail.dispatches. Для SMS отдельного отложенного
топика в текущей реализации нет.

\subsection{Поток статусных событий}

После обработки доставки sender-сервисы публикуют статусные сообщения через
собственные outbox-журналы:
\begin{enumerate}
    \item Mail Sender формирует MailDeliveryStatusChanged и публикует его в notification.mail.delivery-statuses;
    \item SMS Sender формирует SmsDeliveryStatusChanged и публикует его в notification.sms.delivery-statuses;
    \item History Writer читает оба топика и нормализует запись в delivery\_history.
\end{enumerate}

За счёт такой схемы сервис истории полностью отделён от sender-логики и не
вмешивается в принятие решения о доставке. Он работает как отдельный
асинхронный потребитель фактов доставки.

\subsection{Архитектурные свойства Kafka-слоя}

Использование Kafka в данной системе обеспечивает:
\begin{enumerate}
    \item разделение времени приёма команды и времени фактической канальной доставки;
    \item возможность повторной обработки без глобальной транзакции;
    \item масштабирование producer- и consumer-частей независимо друг от друга;
    \item явную наблюдаемость через producer- и listener-метрики Micrometer.
\end{enumerate}

При этом Kafka не заменяет прикладную идемпотентность. Модель доставки системы
остаётся at-least-once, поэтому защита от дублей реализуется в прикладном слое
через уникальные ограничения и входящие журналы \cite{kafka_docs,kafka_design_semantics}.


